% Original English statements and concise English solutions; compile using XeLaTeX.
\documentclass[UTF8,10pt,letterpaper,fontset=fandol]{ctexart}
\usepackage[margin=1in]{geometry}
\usepackage{amsmath,amssymb}
\usepackage{enumitem}
\usepackage{needspace}
\usepackage[unicode,hidelinks]{hyperref}
\pagestyle{plain}
\raggedbottom
\setlength{\parindent}{0pt}
\setlength{\parskip}{0.2em}
\setlength{\emergencystretch}{2em}
\setlist[enumerate]{label=(\arabic*),leftmargin=2em,itemsep=0.4em,topsep=0.3em}
\setcounter{secnumdepth}{0}
\ctexset{section={format=\Large\bfseries,beforeskip=1.2ex,afterskip=0.8ex}}
\newcommand{\R}{\mathbb R}
\newcommand{\C}{\mathbb C}
\newcommand{\Span}{\operatorname{span}}
\newcommand{\Solution}{\par\Needspace{5\baselineskip}\textbf{Solution}\par}

\begin{document}
\begin{center}
{\LARGE Homework 1}\par
\vspace{0.6em}
September 30, 2026\par
Name: YuanHao Tang\quad Student ID: 3120265621
\end{center}
\begin{center}
\emph{Some of the problems are collected from the examples in the book.}
\end{center}

\section{Problem 1}
Let \(\mathcal C_{n-1}(t)\) denote the set of all continuous polynomial functions
defined on \([0,n-1]\) with constant slope in \([i-1,i]\),
\(i=1,2,\ldots,n-1\), \(n\) is an integer. With a vector
\(\alpha(t)\in\mathcal C_{n-1}(t)\), we associate an \(n\)-dimensional
column vector \(\alpha=[\alpha(0)\ \alpha(1)\ \ldots\ \alpha(n-1)]^T\).

\begin{enumerate}
\item Show that the vectors \(\alpha_i(t)\in\mathcal C_{n-1}(t)\),
\(i=1,2,\ldots,n\), are linearly independent if and only if so are the
associated column vectors \(\alpha_i\in\R^n\), \(i=1,2,\ldots,n\).
\item Show that \(\alpha_i(t)\in\mathcal C_3(t)\), \(i=1,2,3,4\) with
\[
\alpha_1=\begin{bmatrix}1\\0\\0\\0\end{bmatrix},\quad
\alpha_2=\begin{bmatrix}-1\\1\\0\\0\end{bmatrix},\quad
\alpha_3=\begin{bmatrix}0\\-1\\1\\0\end{bmatrix},\quad
\alpha_4=\begin{bmatrix}0\\0\\-1\\1\end{bmatrix}
\]
are linearly independent thus form a basis for \(\mathcal C_3(t)\) and so do
the vectors \(\beta_i(t)\in\mathcal C_3(t)\), \(i=1,2,3,4\) with
\[
\beta_1=\begin{bmatrix}1\\0\\0\\0\end{bmatrix},\quad
\beta_2=\begin{bmatrix}1\\1\\0\\0\end{bmatrix},\quad
\beta_3=\begin{bmatrix}1\\1\\1\\0\end{bmatrix},\quad
\beta_4=\begin{bmatrix}1\\1\\1\\1\end{bmatrix}.
\]
\item Find the transition matrix \(T\) that transforms the basis
\([\alpha_1(t)\ \alpha_2(t)\ \alpha_3(t)\ \alpha_4(t)]\) to the basis
\([\beta_1(t)\ \beta_2(t)\ \beta_3(t)\ \beta_4(t)]\).
\item Find the representation of \(\xi(t)\) with
\[
\xi=[\xi(0)\ \xi(1)\ \xi(2)\ \xi(3)]^T
\]
with respect to the basis \([\beta_1(t)\ \beta_2(t)\ \beta_3(t)\ \beta_4(t)]\).
\end{enumerate}

\Solution
\begin{enumerate}
\item The nodal map \(E(f)=[f(0),\ldots,f(n-1)]^T\) is linear.
Any nodal values \(v_0,\ldots,v_{n-1}\) uniquely determine each segment:
\[
f(t)=(i-t)v_{i-1}+(t-i+1)v_i,\qquad t\in[i-1,i].
\]
Thus \(E\) is a linear isomorphism, \(\dim\mathcal C_{n-1}(t)=n\), and
\[
\sum_{i=1}^n c_i\alpha_i(t)=0
\quad\Longleftrightarrow\quad
\sum_{i=1}^n c_i\alpha_i=0.
\]
This proves the equivalence of linear independence.
\item Let \(A_\alpha=[\alpha_1\ \alpha_2\ \alpha_3\ \alpha_4]\) and
\(A_\beta=[\beta_1\ \beta_2\ \beta_3\ \beta_4]\).
Both are upper triangular with unit diagonal, so both determinants are one.
By (1) and \(\dim\mathcal C_3(t)=4\), both lists of functions are bases.
\item From \([\beta_1(t)\ \cdots\ \beta_4(t)]
=[\alpha_1(t)\ \cdots\ \alpha_4(t)]T\), we obtain
\[
\boxed{T=A_\alpha^{-1}A_\beta=
\begin{bmatrix}
1&2&3&4\\
0&1&2&3\\
0&0&1&2\\
0&0&0&1
\end{bmatrix}}.
\]
\item If \(y\) is the required coordinate column, then \(\xi=A_\beta y\). Hence
\[
\boxed{y=A_\beta^{-1}\xi=
\begin{bmatrix}
\xi(0)-\xi(1)\\
\xi(1)-\xi(2)\\
\xi(2)-\xi(3)\\
\xi(3)
\end{bmatrix}}.
\]
\end{enumerate}

\section{Problem 2}
Please show that the vectors
\[
A_1=\begin{bmatrix}0&1\\1&1\end{bmatrix},\quad
A_2=\begin{bmatrix}1&1\\0&1\end{bmatrix},\quad
A_3=\begin{bmatrix}1&0\\1&1\end{bmatrix},\quad
A_4=\begin{bmatrix}1&1\\1&0\end{bmatrix}
\]
are linearly independent and form a basis for \(\C^{2\times2}\), and find the
presentation for the given matrix
\[
A=\begin{bmatrix}1&0\\0&1\end{bmatrix}
\]
with respect to this basis.

\Solution
Suppose \(\sum_{i=1}^4 c_iA_i=0\) and set \(s=\sum_{i=1}^4 c_i\).
Comparing entries gives \(s-c_i=0\) for \(i=1,2,3,4\), so every \(c_i=s\).
Then \(s=4s\) forces \(s=0\), proving linear independence.
Since \(\dim_{\C}\C^{2\times2}=4\), these matrices form a basis.
\[
A=\frac{-A_1+2A_2+2A_3-A_4}{3},\qquad
\boxed{[A]_{(A_1,A_2,A_3,A_4)}=
\begin{bmatrix}-\tfrac13&\tfrac23&\tfrac23&-\tfrac13\end{bmatrix}^{T}}.
\]

\section{Problem 3}
Let \(\mathcal C_{n-1}(t)\) denote the set of all continuous polynomial
functions defined on \([0,n-1]\) with constant slope in \([i-1,i]\),
\(i=1,2,\ldots,n-1\), \(n\) is an integer.
Let \(\boldsymbol\alpha_1,\boldsymbol\alpha_2,\ldots,\boldsymbol\alpha_s\)
be a set of vectors of \(\mathcal V\), the linear span of
\(\boldsymbol\alpha_1,\boldsymbol\alpha_2,\ldots,\boldsymbol\alpha_s\)
is referred to as the set of all possible linear representations of these vectors
\[
\Span\{\boldsymbol\alpha_1,\boldsymbol\alpha_2,\ldots,\boldsymbol\alpha_s\}
=\{\boldsymbol\alpha:\boldsymbol\alpha=k_1\boldsymbol\alpha_1
+k_2\boldsymbol\alpha_2+\cdots+k_s\boldsymbol\alpha_s,\ k_i\in\mathcal F\}.
\]
In this problem, let \(\alpha_1(t)\) and \(\alpha_2(t)\in\mathcal C_3(t)\) with
\[
\alpha_1=\begin{bmatrix}1\\0\\0\\1\end{bmatrix},\quad
\alpha_2=\begin{bmatrix}1\\1\\0\\0\end{bmatrix},\qquad
\mathcal U=\Span\{\alpha_1(t),\alpha_2(t)\},
\]
\(\beta_1(t)\) and \(\beta_2(t)\in\mathcal C_3(t)\) with
\[
\beta_1=\begin{bmatrix}0\\1\\1\\0\end{bmatrix},\quad
\beta_2=\begin{bmatrix}0\\0\\1\\1\end{bmatrix},\qquad
\mathcal V=\Span\{\beta_1(t),\beta_2(t)\}.
\]
\begin{enumerate}
\item The sum of \(\mathcal U\) and \(\mathcal V\) is defined as
\[
\mathcal U+\mathcal V
=\{\alpha=\alpha_1+\alpha_2:\alpha_1\in\mathcal U
\text{ and }\alpha_2\in\mathcal V\}.
\]
Please find out a basis for \(\mathcal U+\mathcal V\) and its dimension;
\item The intersect of \(\mathcal U\) and \(\mathcal V\) is defined as
\[
\mathcal U\cap\mathcal V
=\{\alpha:\alpha\in\mathcal U\text{ and }\alpha\in\mathcal V\}.
\]
Please find out a basis for \(\mathcal U\cap\mathcal V\) and its dimension.
\end{enumerate}

\Solution
By Problem 1, we may work with nodal vectors.
\begin{enumerate}
\item We have \(\beta_2=\alpha_1-\alpha_2+\beta_1\). Moreover,
\[
a\alpha_1+b\alpha_2+c\beta_1
=(a+b,b+c,c,a)^T=0\quad\Longrightarrow\quad a=b=c=0.
\]
Thus a basis and the dimension are
\[
\boxed{\{\alpha_1(t),\alpha_2(t),\beta_1(t)\}},\qquad
\boxed{\dim(\mathcal U+\mathcal V)=3}.
\]
\item Comparing components in \(a\alpha_1+b\alpha_2=c\beta_1+d\beta_2\)
gives \(a=d,\ b=c=-a\). The common vector is \(a(\alpha_1-\alpha_2)\), so
\[
\boxed{\{\alpha_1(t)-\alpha_2(t)\}},\qquad
\boxed{\dim(\mathcal U\cap\mathcal V)=1}.
\]
\end{enumerate}

\section{Examples from the Textbook}
\emph{Read the following examples before attempting Problem 4.}

\textbf{Example 1.23} (p. 21).
Let \(\mathcal D:\R_{n+1}[\lambda]\to\R_n[\lambda]\) be defined as
\[
\mathcal D(f(\lambda))=\frac{d}{d\lambda}f(\lambda)
\qquad\forall f(\lambda)\in\R_{n+1}[\lambda].
\]
Then from the linearity of derivative, it is readily to check that
\(\mathcal D\) is a linear transformation.

\textbf{Example 1.24} (p. 22).
Let \(\mathcal S:\R_n[\lambda]\to\R_{n+1}[\lambda]\) be defined as
\[
\mathcal S(f(\lambda))=\int_0^\lambda f(\xi)\,d\xi
\qquad\forall f(\lambda)\in\R_n[\lambda].
\]
Then from the linearity of integration, it is readily to check that
\(\mathcal S\) is a linear transformation.

\begin{samepage}
\textbf{Example 1.26} (p. 25).
Determine the matrix representation of the linear transformation
\(\mathcal D\in\mathcal L(\R_n[\lambda]\to\R_{n-1}[\lambda])\),
\[
\mathcal D(f(\lambda))=\frac{d}{d\lambda}f(\lambda)
\]
with respect to the bases \([1\ \lambda\ \cdots\ \lambda^n]\) and
\([1\ \lambda\ \cdots\ \lambda^{n-1}]\).
\end{samepage}

\textbf{Example 1.27} (p. 26).
Define the linear transformation
\(\mathcal S\in\mathcal L(\R_n[\lambda]\to\R_{n+1}[\lambda])\) as
\[
\mathcal S(f(\lambda))=\int_0^\lambda f(\xi)\,d\xi
\qquad\forall f(\lambda)\in\R_n[\lambda].
\]
Determine the matrix representation with respect to the basis
\([1\ \lambda\ \cdots\ \lambda^n]\) and \([1\ \lambda\ \cdots\ \lambda^{n+1}]\).

\section{Problem 4}
Denote by \(\R_{n-1}(\lambda)\) the set of all polynomials with real coefficients
and degrees not exceeding \(n-1\).
With respect to the following addition and scalar product operations,
\[
f_1(\lambda)=\sum_{i=0}^{n-1}f_i^{(1)}\lambda^i,\qquad
f_2(\lambda)=\sum_{i=0}^{n-1}f_i^{(2)}\lambda^i
\]
\[
(f_1+f_2)(\lambda)=\sum_{i=0}^{n-1}
\bigl(f_i^{(1)}+f_i^{(2)}\bigr)\lambda^i
\]
\[
(kf_1)(\lambda)=k\sum_{i=0}^{n-1}f_i^{(1)}\lambda^i
\]
\begin{enumerate}
\item Let a linear transformation
\(\mathcal L_\alpha:\R_{n-1}(\lambda)\to\R_{n-1}(\lambda)\) be defined as
\[
\mathcal L_\alpha[f(\lambda+\alpha)]=f(\lambda),\qquad
f(\lambda+\alpha),\ f(\lambda)\in\R_{n-1}(\lambda),
\]
where \(\alpha\) is real constant number, show that
\(\mathcal L_\alpha\) is a linear transformation;
\item Find out the matrix representation \(L\) for the linear transformation
\(\mathcal L_\alpha\) with respect to the basis
\(1,\lambda+\alpha,(\lambda+\alpha)^2,\ldots,(\lambda+\alpha)^{n-1}\),
i.e., the matrix \(L\) satisfying
\[
\begin{aligned}
&\mathcal L_\alpha\bigl([1,\lambda+\alpha,(\lambda+\alpha)^2,\ldots,
(\lambda+\alpha)^{n-1}]\bigr)\\
&\qquad\mathrel{\triangleq}
[1,\lambda+\alpha,(\lambda+\alpha)^2,\ldots,(\lambda+\alpha)^{n-1}]L.
\end{aligned}
\]
\end{enumerate}

\Solution
\begin{enumerate}
\item Set \(g(\lambda)=f(\lambda+\alpha)\). Then
\((\mathcal L_\alpha g)(\lambda)=g(\lambda-\alpha)\), and shifting does not increase degree.
For any \(a,b\in\R\),
\[
\begin{aligned}
\bigl(\mathcal L_\alpha(ag_1+bg_2)\bigr)(\lambda)
&=ag_1(\lambda-\alpha)+bg_2(\lambda-\alpha)\\
&=a(\mathcal L_\alpha g_1)(\lambda)+b(\mathcal L_\alpha g_2)(\lambda).
\end{aligned}
\]
Thus \(\mathcal L_\alpha\) is linear.
\item Let \(b_j(\lambda)=(\lambda+\alpha)^j\), \(0\le j\le n-1\). Then
\[
\mathcal L_\alpha b_j=\lambda^j
=\sum_{i=0}^j\binom ji(-\alpha)^{j-i}b_i.
\]
Placing these coordinates in columns gives
\[
\boxed{L_{i+1,j+1}=\begin{cases}
\displaystyle\binom ji(-\alpha)^{j-i},&0\le i\le j\le n-1,\\
0,&0\le j<i\le n-1.
\end{cases}}
\]
\end{enumerate}

\end{document}

